% ============================================================
%  Sprawozdanie: Badanie polaryzacji światła (02_POLAR)
%  Fotonika – laboratorium
% ============================================================
\documentclass[a4paper,12pt]{article}

\usepackage[utf8]{inputenc}
\usepackage[T1]{fontenc}
\usepackage[polish]{babel}
\usepackage{geometry}
\geometry{a4paper, top=2.2cm, bottom=2.2cm, left=2.5cm, right=2.5cm}

\usepackage{amsmath,amssymb}
\usepackage{booktabs}
\usepackage{graphicx}
\usepackage{float}
\usepackage{caption}
\usepackage{multirow}
\usepackage{array}
\usepackage{xcolor}
\usepackage{hyperref}
\usepackage{siunitx}
\sisetup{output-decimal-marker={,}, locale=DE}
\usepackage{parskip}
\setlength{\parskip}{4pt}

\title{}
\date{}

\begin{document}

% ============================================================
%  TABELA NAGŁÓWKOWA -- do samodzielnego uzupełnienia
% ============================================================

% ============================================================
%  SCHEMAT UKŁADU POMIAROWEGO -- do samodzielnego uzupełnienia
% ============================================================

% ============================================================
\section*{4.\quad Cel ćwiczenia}
% ============================================================

Celem ćwiczenia jest zbadanie zjawiska polaryzacji światła laserowego.
W~ramach ćwiczenia wyznaczono kąt Brewstera dla płytki kwarcowej, a~na
jego podstawie obliczono współczynnik załamania materiału. Następnie
zbadano właściwości polaryzacyjne stosu polaryzacyjnego (15~płytek szklanych)
i~wyznaczono stopień polaryzacji wiązki przepuszczonej przez stos oraz
dynamikę polaryzatora.

% ============================================================
\section*{5.\quad Wstęp teoretyczny}
% ============================================================

\textbf{Polaryzacja światła.}
Światło jest poprzeczną falą elektromagnetyczną; jego stan polaryzacji opisuje
sposób, w~jaki wektor pola elektrycznego $\vec{E}$ porusza się w~czasie.
W~świetle \emph{niespolaryzowanym} kierunek drgań zmienia się chaotycznie;
w~świetle \emph{spolaryzowanym liniowo} wektor $\vec{E}$ drga w~jednej
stałej płaszczyźnie zawierającej kierunek propagacji.

\medskip
\textbf{Prawo Malusa.}
Gdy wiązka spolaryzowana liniowo o~natężeniu $I_0$ przechodzi przez analizator
ustawiony pod kątem $\varphi$ względem płaszczyzny polaryzacji, natężenie
wiązki wychodzące wynosi:
\begin{equation}
  I = I_0 \cos^2\!\varphi .
  \label{eq:malus}
\end{equation}
Przy \(\varphi = 90^\circ\) wiązka zostaje całkowicie wygaszona.

\medskip
\textbf{Kąt Brewstera i~prawo Brewstera.}
Gdy wiązka padająca na granicę ośrodków pod kątem Brewstera $\theta_B$ jest
spolaryzowana w~płaszczyźnie padania (składowa~p), składowa~p nie ulega
odbiciu ($r_p = 0$), a~wiązka odbita zawiera wyłącznie składową s
(polaryzacja pionowa). Warunek ten zachodzi, gdy wiązka odbita jest
prostopadła do~wiązki załamanej, co daje prawo Brewstera:
\begin{equation}
  n = \tan\theta_B ,
  \label{eq:brewster}
\end{equation}
gdzie $n$ -- współczynnik załamania ośrodka, $n_1 = 1$ (powietrze).

\medskip
\textbf{Stos polaryzacyjny.}
Stos polaryzacyjny zbudowany jest z~$N$ płaskorównoległych płytek szklanych
ustawionych pod kątem Brewstera. Przy każdym przejściu przez powierzchnię
płytki składowa~p przechodzi bez straty, a~część składowej~s ulega odbiciu.
Natężeniowy współczynnik odbicia składowej~s przy kącie Brewstera:
$R_s = \bigl(\frac{1-n^2}{1+n^2}\bigr)^2$.
Po przejściu przez $N$ płytek ($2N$ powierzchni):
\begin{equation}
  I_s^{(\text{out})} = I_s^{(\text{in})}\cdot(1 - R_s)^{2N} .
  \label{eq:stos}
\end{equation}

\textbf{Stopień polaryzacji} wiązki wychodzące ze stosu wyraża się wzorem:
\begin{equation}
  P = \frac{I_{\max} - I_{\min}}{I_{\max} + I_{\min}} ,
  \label{eq:P}
\end{equation}
gdzie $I_{\max}$ i~$I_{\min}$ są odpowiednio maksymalną i~minimalną mocą
optyczną rejestrowaną przy obrocie analizatora.

\textbf{Dynamika polaryzatora} (stosunek wygaszenia) wyrażona w~decybelach:
\begin{equation}
  D = 10\log_{10}\frac{I_{\max}}{I_{\min}} .
  \label{eq:dB}
\end{equation}

% ============================================================
\section*{7.\quad Spis przyrządów pomiarowych}
% ============================================================

\begin{itemize}
  \item Laser He-Ne ($\lambda = 632{,}8\ \si{nm}$) -- źródło spójnego,
        spolaryzowanego liniowo światła.
  \item Obrotowy polaryzator liniowy (P) -- ustawienie osi transmisji
        w~płaszczyźnie poziomej (polaryzacja~p).
  \item Płytka kwarcowa na stoliku obrotowym ze skalą kątową
        (dokładność odczytu: \(1^\circ\)).
  \item Stos polaryzacyjny -- 15~płytek szklanych ustawionych pod kątem
        Brewstera.
  \item Analizator (A) -- obrotowy polaryzator liniowy za stosem.
  \item Fotodetektor (FD) z~miernikiem mocy optycznej.
  \item Ekran do obserwacji wzrokowej plamki świetlnej.
\end{itemize}

% ============================================================
\section*{8.\quad Przebieg ćwiczenia}
% ============================================================

\textbf{Część~A -- Stan polaryzacji wiązki laserowej.}
Zestawiono układ: Laser $\to$ Polaryzator~P $\to$ Fotodetektor~FD.
Obracano polaryzator~P o~pełny obrót (\(360^\circ\)) i~obserwowano zmiany mocy
na~mierniku. Wygaszenie wiązki przy pewnym ustawieniu~P potwierdziło, że
laser emituje światło liniowo spolaryzowane.

\medskip
\textbf{Część~B -- Wyznaczenie kąta Brewstera.}
Do układu włączono płytkę kwarcową na stoliku obrotowym, ustawioną pod kątem
prostym do wiązki (pozycja zerowa). Polaryzator~P ustawiono tak, aby oś
transmisji była pozioma (polaryzacja~p). Wolno obracano stolik, obserwując
plamkę odbitą, aż do jej całkowitego zaniku. Odczytano wskazanie skali
kątowej jako kąt Brewstera $\theta_B$.

\medskip
\textbf{Część~C -- Badanie stosu polaryzacyjnego.}
Płytkę kwarcową zastąpiono stosem 15~płytek szklanych ustawionym pod
kątem Brewstera. Za stosem umieszczono analizator~A i~fotodetektor~FD.
Obracano analizator~A o~\(360^\circ\), rejestrując wartości maksymalnej
($I_{\max}$) i~minimalnej ($I_{\min}$) mocy optycznej.

% ============================================================
\section*{9.\quad Wyniki pomiarów i~opracowanie}
% ============================================================

% ─────────────────────────────────────────────────────────────
\subsection*{9.1\quad Część~B -- Kąt Brewstera i~współczynnik załamania}

Zmierzony kąt Brewstera:
\[
  \theta_B = 56^\circ.
\]

Współczynnik załamania obliczony z~prawa Brewstera~(\ref{eq:brewster}):
\[
  n = \tan 56^\circ = 1{,}4826.
\]

\noindent\textbf{Analiza niepewności wyznaczenia $n$:}

Niepewność instrumentalna kąta (typ~B) wynikająca z~podziałki stolika
(dokładność odczytu $\Delta\theta_B = 1^\circ$):
\[
  u_B(\theta_B) = \frac{\Delta\theta_B}{\sqrt{3}}
  = \frac{1^\circ}{\sqrt{3}}
  \approx 0{,}577^\circ
  = 0{,}01007\ \si{rad}.
\]

Niepewność współczynnika załamania metodą różniczki:
\[
  \Delta n
  = \frac{u_B(\theta_B)}{\cos^2\!\theta_B}
  = \frac{0{,}01007}{\cos^2\!56^\circ}
  = \frac{0{,}01007}{0{,}3127}
  \approx 0{,}032.
\]

Wynik końcowy (z~niepewnością rozszerzoną $k = 2$):
\[
  \boxed{n = 1{,}483 \pm 0{,}064}
\]

\noindent Dla porównania: tabelaryczny współczynnik załamania kwarcu dla
$\lambda = 632{,}8\ \si{nm}$ wynosi $n_{\text{lit}} \approx 1{,}457$.
Zmierzona wartość jest zbliżona do wartości literaturowej; odchylenie wynosi
około $1{,}8\%$, co mieści się w~granicy wyznaczonej niepewności.

% ─────────────────────────────────────────────────────────────
\subsection*{9.2\quad Część~C -- Stopień polaryzacji stosu}

\noindent Zmierzone moce optyczne za analizatorem:
\begin{align*}
  I_{\max} &= I_{\text{pozioma}} = 0{,}38\ [\text{j.u.}],\\
  I_{\min} &= I_{\text{pionowa}}  = 0{,}03\ [\text{j.u.}].
\end{align*}

Kierunek $I_{\max}$ odpowiada osi poziomej analizatora, a~$I_{\min}$ osi
pionowej -- jest to zgodne z~oczekiwaniem: stos przepuszcza głównie
składową~p (poziomą), a~składową~s (pionową) odbija.

\medskip
\noindent\textbf{Stopień polaryzacji:}
\[
  P = \frac{I_{\max} - I_{\min}}{I_{\max} + I_{\min}}
    = \frac{0{,}38 - 0{,}03}{0{,}38 + 0{,}03}
    = \frac{0{,}35}{0{,}41}
    \approx 0{,}854.
\]

\noindent\textbf{Dynamika polaryzatora (stosunek wygaszenia):}
\[
  D = 10\log_{10}\frac{I_{\max}}{I_{\min}}
  = 10\log_{10}\frac{0{,}38}{0{,}03}
  \approx 10 \cdot 1{,}103
  = 11{,}03\ \si{dB}.
\]

Zmierzona (oszacowana) dynamika polaryzatora wynosi $D = 10{,}414\ \si{dB}$,
co odpowiada stosunkowi $I_{\max}/I_{\min} \approx 11{,}0$. Niewielka
różnica względem obliczonego ilorazu $0{,}38/0{,}03 = 12{,}67$ wynika
z~zaokrąglenia zarejestrowanych wskazań miernika.

\medskip
\noindent\textbf{Analiza niepewności stopnia polaryzacji~$P$:}

Zakładając rozdzielczość miernika mocy $\Delta I = 0{,}005\ [\text{j.u.}]$
(niepewność instrumentalna, typ~B):
\[
  u(I_{\max}) = u(I_{\min}) = \frac{\Delta I}{\sqrt{3}}
  \approx \frac{0{,}005}{1{,}732}
  \approx 0{,}0029\ [\text{j.u.}].
\]

Propagacja niepewności dla $P = (I_{\max} - I_{\min})/(I_{\max} + I_{\min})$:
\begin{align*}
  \frac{\partial P}{\partial I_{\max}}
    &= \frac{2I_{\min}}{(I_{\max}+I_{\min})^2}
     = \frac{2 \cdot 0{,}03}{(0{,}41)^2}
     = \frac{0{,}06}{0{,}1681}
     \approx 0{,}357\ [\text{j.u.}]^{-1}, \\
  \frac{\partial P}{\partial I_{\min}}
    &= \frac{-2I_{\max}}{(I_{\max}+I_{\min})^2}
     = \frac{-2 \cdot 0{,}38}{0{,}1681}
     \approx -4{,}521\ [\text{j.u.}]^{-1}.
\end{align*}

\[
  u(P) = \sqrt{\left(\frac{\partial P}{\partial I_{\max}}\right)^{\!2}
               u^2(I_{\max})
              +\left(\frac{\partial P}{\partial I_{\min}}\right)^{\!2}
               u^2(I_{\min})}
       = \sqrt{(0{,}357)^2(0{,}0029)^2 + (4{,}521)^2(0{,}0029)^2}
       \approx 0{,}013.
\]

Wynik z~niepewnością rozszerzoną ($k = 2$):
\[
  \boxed{P = 0{,}854 \pm 0{,}027}
\]

% ─────────────────────────────────────────────────────────────
\subsection*{9.3\quad Podsumowanie wyników}

\begin{table}[H]
\centering
\caption{Zestawienie wyników pomiarów.}
\label{tab:wyniki}
\renewcommand{\arraystretch}{1.3}
\begin{tabular}{lll}
\toprule
Wielkość & Wartość & Uwagi \\
\midrule
Kąt Brewstera $\theta_B$ & $56^\circ$ & odczyt ze skali stolika \\
Współczynnik załamania $n$ & $1{,}483 \pm 0{,}064$ & prawo Brewstera \\
Wartość literaturowa $n_{\text{lit}}$ & $\approx 1{,}457$ & kwarc, $\lambda = 632{,}8\ \si{nm}$ \\
\midrule
Moc maks.\ za analizatorem $I_{\max}$ & $0{,}38\ [\text{j.u.}]$ & pol.\ pozioma \\
Moc min.\ za analizatorem $I_{\min}$ & $0{,}03\ [\text{j.u.}]$ & pol.\ pionowa \\
Stopień polaryzacji $P$ & $0{,}854 \pm 0{,}027$ & wzór~(\ref{eq:P}) \\
Dynamika polaryzatora $D$ & $10{,}414\ \si{dB}$ & pomiar/oszacowanie \\
\bottomrule
\end{tabular}
\end{table}

% ============================================================
\section*{10.\quad Wnioski}
% ============================================================

\begin{enumerate}

  \item \textbf{Laser He-Ne emituje światło spolaryzowane liniowo.}
        Przy obrocie polaryzatora o~$90^\circ$ wiązka laserowa zostaje całkowicie
        wygaszona, co jest charakterystyczne dla polaryzacji liniowej
        (zgodnie z~prawem Malusa).

  \item \textbf{Wyznaczony kąt Brewstera} dla płytki kwarcowej wynosi
        $\theta_B = 56^\circ$, co daje współczynnik załamania
        $n = \tan 56^\circ = 1{,}483 \pm 0{,}064$.
        Wartość ta jest zgodna z~tabelaryczną wartością współczynnika
        załamania kwarcu $n_{\text{lit}} \approx 1{,}457$
        (odchylenie $\approx 1{,}8\%$), mieszcząc się w~granicach
        wyznaczonej niepewności. Dominującym źródłem niepewności jest
        rozdzielczość skali kątowej stolika obrotowego (\(1^\circ\)).

  \item \textbf{Odbite od płytki kwarcowej pod kątem Brewstera}
        światło jest spolaryzowane pionowo (wyłącznie składowa~s),
        co potwierdzono obserwacją zaniku wiązki odbitej przy obrocie
        analizatora do orientacji poziomej.

  \item \textbf{Stos polaryzacyjny skutecznie polaryzuje wiązkę przepuszczoną.}
        Zmierzone moce $I_{\max} = 0{,}38\ \text{j.u.}$ (polaryzacja pozioma)
        i~$I_{\min} = 0{,}03\ \text{j.u.}$ (polaryzacja pionowa) dają stopień
        polaryzacji $P = 0{,}854 \pm 0{,}027 \approx 85\%$.
        Wartość ta jest nieco niższa od teoretycznej dla 15~płytek ($P \approx 98\%$),
        co może wynikać ze~skończonej jakości ustawienia kąta
        Brewstera lub strat w~ośrodku.

  \item \textbf{Kierunek polaryzacji} wiązki przepuszczonej przez stos jest
        \emph{poziomy} (składowa~p), zaś wiązka odbita od stosu jest
        spolaryzowana \emph{pionowo} (składowa~s). Wynik ten jest zgodny z~teorią:
        stos selektywnie pochłania składową~s, transmitując składową~p praktycznie
        bez strat.

  \item \textbf{Dynamika polaryzatora} wynosi $D = 10{,}414\ \si{dB}$
        ($I_{\max}/I_{\min} \approx 11{,}0$), co wskazuje, że stos tłumi
        niepożądaną składową polaryzacji o~ponad rząd wielkości.

\end{enumerate}

\end{document}
